\chapter{Repeated Games}\label{ch:b2:repeated-games}

Firms that agree to fix prices gain together, and each gains more by secretly undercutting the others. In a
game played once, the agreement collapses: it is a prisoner's dilemma. Yet cartels exist, and some last for
decades. The chart summarises how long cartels lasted in nine samples, from nineteenth-century
international cartels to those prosecuted in the 1990s.

\begin{openingchart}{Duration of cartels in nine samples: average in years (dot) and, where reported, the
  shortest and longest duration (bar). For export associations of 1918--1965 the figure is a median;
  for international cartels of 1920--1939 it covers the 28 cartels not ended by the Second World War (8.3 years
  for all 71). Samples sorted by average duration.}{Levenstein and Suslow (2006), Table~1}
\begin{tikzpicture}
\begin{axis}[omrecord, width=0.62\linewidth, height=0.6\linewidth, xmin=0, xmax=30, ymin=0.4,
    ymax=9.6, ymajorgrids=false, xlabel={Duration of a cartel (years)},
    ytick=data, yticklabels from table={figdata/bachelor-2/13-repeated-games/out/cartels.csv}{sample},
    y tick label style={font=\footnotesize}]
  \addplot[only marks, mark=*, mark size=2.2pt, omDef, error bars/.cd, x dir=both, x explicit,
      error bar style={draw=omDef, line width=0.9pt}, error mark options={draw=omDef, rotate=90, mark size=2pt}]
    table[x=mean, y=pos, x error plus=plus, x error minus=minus]
    {figdata/bachelor-2/13-repeated-games/out/cartels.csv};
\end{axis}
\end{tikzpicture}
\end{openingchart}

% ledger: b2:ls2006-t1, b2:ls2006-text
% figdata: cartels_summary (lowest 3.7, highest 10, median 5.3, longest 29, with_range 6)
Average durations run from 3.7 years (international cartels of 1920--1939, counting only those not ended by the war) to
10 years, with a median of 5.3, and within each sample the spread is wide:
some cartels end within a year, others last 18 to 29 years. These are cartels that were discovered, most of
them by prosecutors, so the sample leaves out any that were never found.

The prisoner's dilemma of \cref{ch:b2:static-games} is played once. Firms in a market meet again every week,
and what they do today can change what the others do tomorrow. This chapter shows that repetition can sustain
cooperation that a single play cannot, works out when, and draws the lessons for \omterm{def:b2:repeated-games:collusion}{collusion} and for
\omterm{def:b2:repeated-games:reputation}{reputations} in markets where buyers cannot see quality before they buy.

\begin{recall}
From \cref{ch:b2:static-games}: the prisoner's dilemma and \omterm{def:b2:static-games:nash-equilibrium}{Nash equilibrium}. From
\cref{ch:b2:dynamic-games}: \omterm{def:b2:dynamic-games:spe}{subgame-perfect equilibrium} and \omterm{def:b2:dynamic-games:backward-induction}{backward induction}. From the Year 1 volume:
Bertrand and Cournot competition.
\end{recall}

\begin{notation}[Symbols of this chapter]\label{not:b2:repeated-games:symbols}
In the prisoner's dilemma, each player earns $R$ when both cooperate, $P$ when both defect, $T$ when she defects
against a cooperator and $S$ when she cooperates against a defector, with $T > R > P > S$. $\beta$ is the \omterm{def:b2:intertemporal-choice:discount-factor}{discount factor}
per period, $\beta^*$ its critical value, $n$ the number of firms, $\Pi$ the monopoly profit of the market per period and $L$
the number of periods before a deviation is detected.
\end{notation}

\section{Finite repetition}\label{sec:b2:repeated-games:finite}

\begin{definition}[Stage game]\label{def:b2:repeated-games:stage-game}
The \emph{stage game}\index{stage game} of a \omterm{def:b2:repeated-games:repeated-game}{repeated game} is the game in normal form that is played in each
period.
\end{definition}

\begin{definition}[Repeated game]\label{def:b2:repeated-games:repeated-game}
A \emph{repeated game}\index{repeated game} is a dynamic game in which the same \omterm{def:b2:repeated-games:stage-game}{stage game} is played in
every period, finitely or infinitely often, each player observing the actions of every earlier period before
choosing her action in the current one.
\end{definition}

A strategy in a \omterm{def:b2:repeated-games:repeated-game}{repeated game} says what to play in each period after each possible history of play.

\begin{proposition}[Unravelling]\label{prop:b2:repeated-games:unravelling}
If the prisoner's dilemma is repeated a known, finite number of times, every player defects in every period in
the unique \omterm{def:b2:dynamic-games:spe}{subgame-perfect equilibrium}.
\end{proposition}

\begin{proof}
In the last period nothing follows, so the period is the \omterm{def:b2:repeated-games:stage-game}{stage game}, and both defect whatever happened before.
In the period before, play in the last period is therefore fixed whatever happens now; today's action has no
effect on the future, so both defect today too. Repeating the argument back to the first period, both defect in
every period after every history.
\end{proof}

The argument relies on a last period that both players know. When the game might go on, the future can
reward cooperation today.

\section{Infinite repetition and discounting}\label{sec:b2:repeated-games:infinite}

\begin{definition}[Discounted average payoff]\label{def:b2:repeated-games:average-payoff}
The \emph{discounted average payoff}\index{discounted average payoff} of a stream of stage-game payoffs $u_0,
u_1, u_2, \dots$ is $(1 - \beta)\sum_{t \ge 0} \beta^t u_t$. A constant stream $u$ has discounted average payoff $u$.
\end{definition}

The \omterm{def:b2:intertemporal-choice:discount-factor}{discount factor} $\beta$ need not measure patience alone. If the game continues to the next period with
probability $\rho$ and the interest rate is $r$, a payoff next period is worth $\rho/(1 + r)$ today: an infinitely \omterm{def:b2:repeated-games:repeated-game}{repeated game}
with $\beta = \rho/(1 + r)$ describes a relationship that ends at a random date that no one can foresee.

\section{Trigger strategies}\label{sec:b2:repeated-games:trigger}

\begin{definition}[Grim-trigger strategy]\label{def:b2:repeated-games:grim-trigger}
In the repeated prisoner's dilemma, the \emph{grim-trigger strategy}\index{grim-trigger strategy} cooperates in
the first period and in every period in which no player has defected before, and defects for ever after any
defection.
\end{definition}

\begin{definition}[Tit for tat]\label{def:b2:repeated-games:tit-for-tat}
\emph{Tit for tat}\index{tit for tat} cooperates in the first period and then plays whatever the other player
played in the previous period.
\end{definition}

\begin{definition}[Critical discount factor]\label{def:b2:repeated-games:critical-delta}
The \emph{critical discount factor}\index{critical discount factor} $\beta^*$ of a \omterm{def:b2:static-games:profile}{strategy profile} is the smallest
\omterm{def:b2:intertemporal-choice:discount-factor}{discount factor} at which the profile is a \omterm{def:b2:dynamic-games:spe}{subgame-perfect equilibrium}.
\end{definition}

To check that a profile is \omterm{def:b2:dynamic-games:subgame}{subgame} perfect, it is enough to check one-shot deviations: that in no \omterm{def:b2:dynamic-games:subgame}{subgame} can a
player gain by changing her action in one period and then following her strategy again. If a deviation of any
kind gained, then with discounting a deviation that ends after finitely many periods would gain too, and the
last period in which such a deviation differs from the strategy would be a profitable one-shot deviation.

\begin{proposition}[Grim trigger]\label{prop:b2:repeated-games:grim}
In the infinitely repeated prisoner's dilemma, both players playing grim trigger is a \omterm{def:b2:dynamic-games:spe}{subgame-perfect
equilibrium} if and only if
\[
\beta \ge \beta^* = \frac{T - R}{T - P}.
\]
\end{proposition}

\begin{proof}
There are two kinds of \omterm{def:b2:dynamic-games:subgame}{subgame}. After a defection, both defect for ever; this repeats a \omterm{def:b2:static-games:nash-equilibrium}{Nash equilibrium} of the
\omterm{def:b2:repeated-games:stage-game}{stage game}, and since future play does not depend on today's action, no one-shot deviation gains. With no
defection so far, cooperating for ever yields $R$; a one-shot deviation yields $T$ today and then $P$ for ever, a
discounted average of $(1 - \beta)T + \beta P$. Cooperation is a \omterm{def:b2:static-games:best-response}{best response} if and only if $R \ge (1 - \beta)T + \beta P$, that is,
$\beta \ge (T - R)/(T - P)$.
\end{proof}

With $T = 5$, $R = 3$, $P = 1$, $S = 0$, the \omterm{def:b2:repeated-games:critical-delta}{critical discount factor} is $2/4 = 0.5$. The listing computes the \omterm{def:b2:repeated-games:average-payoff}{discounted
average payoffs} of four strategies against each other, at a \omterm{def:b2:intertemporal-choice:discount-factor}{discount factor} below the critical value and at one
above it.

\omcode{code/bachelor-2/13-repeated-games/listings/repeated_pd.py}{1}{28}{\omterm{def:b2:repeated-games:average-payoff}{Discounted average payoffs} of four
strategies in the repeated prisoner's dilemma, at two \omterm{def:b2:intertemporal-choice:discount-factor}{discount factors}.}\label{lst:b2:repeated-games:pd}

\omoutput{code/bachelor-2/13-repeated-games/outputs/repeated_pd.txt}{Against grim trigger, always defecting pays
3.4 when $\beta = 0.4$ and only 1.4 when $\beta = 0.9$.}\label{out:b2:repeated-games:pd}

% code: listings/repeated_pd.py
At $\beta = 0.4$ a player facing grim trigger earns 3.4 by defecting for ever against 3 by cooperating: the gain
today outweighs the loss to come. At $\beta = 0.9$ the same deviation earns 1.4. \omterm{def:b2:repeated-games:tit-for-tat}{Tit for tat} sustains cooperation
too, but it punishes for one period only and then forgives, so it needs more patience against some
deviations (\cref{exo:b2:repeated-games:6}).

\section{The folk theorem}\label{sec:b2:repeated-games:folk}

Grim trigger sustains the payoffs $(3, 3)$. Other strategies sustain other payoffs: for instance, alternating
between mutual cooperation and mutual defection, with any defection from the pattern punished for ever. The
range of what repetition can sustain is wide.

\begin{definition}[Folk theorem]\label{def:b2:repeated-games:folk-theorem}
A \emph{folk theorem}\index{folk theorem} is a result stating that, in an infinitely \omterm{def:b2:repeated-games:repeated-game}{repeated game} with patient
enough players, every feasible payoff vector that gives each player more than the least she can guarantee
herself in the \omterm{def:b2:repeated-games:stage-game}{stage game}, whatever the others do, is the \omterm{def:b2:repeated-games:average-payoff}{discounted average payoff} of a \omterm{def:b2:dynamic-games:spe}{subgame-perfect
equilibrium}.
\end{definition}

\begin{theorem}[Folk theorem, with Nash punishment]\label{thm:b2:repeated-games:folk}
Let $e$ be a \omterm{def:b2:static-games:nash-equilibrium}{Nash equilibrium} of the \omterm{def:b2:repeated-games:stage-game}{stage game}, and $v$ a feasible payoff vector with $v_i > u_i(e)$ for every
player. There is $\beta^* < 1$ such that for every $\beta \ge \beta^*$ the \omterm{def:b2:repeated-games:repeated-game}{repeated game} has a \omterm{def:b2:dynamic-games:spe}{subgame-perfect equilibrium} with
\omterm{def:b2:repeated-games:average-payoff}{discounted average payoffs} $v$.
\end{theorem}

\admitted

For the prisoner's dilemma, with $e$ mutual defection and $v = (R, R)$, \cref{prop:b2:repeated-games:grim} is the
proof, with the exact $\beta^*$. The figure shows the payoffs the theorem covers. In the prisoner's dilemma each player
can always guarantee herself 1, by defecting, so no equilibrium gives a player less
(\cref{exo:b2:repeated-games:10}).

\begin{omfigure}
\begin{tikzpicture}[econ]
\begin{axis}[omrecord, grid=none, xmin=-0.3, xmax=6.3, ymin=-0.3, ymax=5.6, width=0.52\linewidth,
    height=0.5\linewidth, xlabel={Player 1's average payoff}, ylabel={Player 2's average payoff},
    xtick={0,1,3,5}, ytick={0,1,3,5}]
  \addplot[fill=omDef!12, draw=black!60] coordinates {(1,1) (5,0) (3,3) (0,5) (1,1)};
  \addplot[fill=omDef!40, draw=none] coordinates {(1,1) (4.3333,1) (3,3) (1,4.3333) (1,1)};
  \draw[guide] (axis cs:1,-0.3) -- (axis cs:1,5.6);
  \draw[guide] (axis cs:-0.3,1) -- (axis cs:6.3,1);
  \node[eqpoint] at (axis cs:3,3) {};
  \node[eqpoint] at (axis cs:1,1) {};
  \node[anchor=south west, font=\footnotesize] at (axis cs:3.05,3.05) {$(R, R)$};
  \node[anchor=north east, font=\footnotesize] at (axis cs:0.95,0.95) {$(P, P)$};
  \node[anchor=west, font=\footnotesize] at (axis cs:5,0.15) {$(T, S)$};
  \node[anchor=south, font=\footnotesize] at (axis cs:0.4,5) {$(S, T)$};
\end{axis}
\end{tikzpicture}

{\small Payoffs in the repeated prisoner's dilemma with $T = 5$, $R = 3$, $P = 1$, $S = 0$. The light region contains
every feasible pair of average payoffs, the averages of the four stage-game outcomes. The dark region holds the
feasible pairs that give both players more than 1, what each can guarantee by defecting; with patient enough
players, each of them is the outcome of a \omterm{def:b2:dynamic-games:spe}{subgame-perfect equilibrium}.}
\end{omfigure}

The \omterm{def:b2:repeated-games:folk-theorem}{folk theorem} is a result about possibility. It says that repetition makes cooperation possible, and also
that almost anything else is possible: the same game has equilibria with full cooperation, with none, and with
everything in between. Which one is played the theory does not say.

\section{When collusion can be sustained}\label{sec:b2:repeated-games:collusion}

\begin{definition}[Collusion]\label{def:b2:repeated-games:collusion}
\emph{Collusion}\index{collusion} is an equilibrium of a \omterm{def:b2:repeated-games:repeated-game}{repeated game} between firms in which prices are held
above, or quantities below, those of the equilibrium of the \omterm{def:b2:repeated-games:stage-game}{stage game}, sustained by the threat that a deviation
will be followed by lower profits. It is \emph{tacit} when it involves no explicit agreement.
\end{definition}

Take $n$ identical firms competing on price, with constant marginal cost: the stage-game equilibrium is price
equal to cost, with zero profit. If all charge the monopoly price they share the monopoly profit, $\Pi/n$ each. A
firm that undercuts slightly takes the whole market for one period, $\Pi$, and then earns zero if the others
revert to the stage-game equilibrium for ever. \omterm{def:b2:repeated-games:collusion}{Collusion} is sustainable when $\Pi/n \ge (1 - \beta)\Pi$:
\[
\beta \ge \beta^* = 1 - \frac{1}{n}.
\]
Two firms need $\beta \ge 0.5$, five need 0.8, ten need 0.9. The more firms, the smaller each one's share of the
collusive profit, and the more tempting the whole market.

\begin{omfigure}
\begin{tikzpicture}[econ]
\begin{axis}[omrecord, xmin=1.5, xmax=10.5, ymin=0.4, ymax=1.02, width=0.62\linewidth,
    height=0.42\linewidth, xlabel={Number of firms $n$}, ylabel={Critical discount factor $\beta^*$},
    xtick={2,...,10}, legend style={at={(0.5,-0.25)}, anchor=north, legend columns=2,
    /tikz/every even column/.append style={column sep=10pt}}]
  \addplot[omDef, thick, mark=*, mark size=1.6pt, samples at={2,...,10}] {1 - 1/x};
  \addplot[omThm, thick, dashed, mark=square*, mark size=1.4pt, samples at={2,...,10}] {sqrt(1 - 1/x)};
  \legend{Detected at once, Detected after one period}
\end{axis}
\end{tikzpicture}

{\small The \omterm{def:b2:repeated-games:critical-delta}{critical discount factor} for \omterm{def:b2:repeated-games:collusion}{collusion} among $n$ firms competing on price. When a deviation is
detected only after $L$ periods, the deviator keeps the whole market for $L + 1$ periods, and the condition becomes
$\beta^{L + 1} \ge 1 - 1/n$; the dashed line is $L = 1$.}
\end{omfigure}

The same reasoning, applied to the other features of a market, gives the model's comparative statics.
\begin{itemize}
  \item \emph{Detection lags.} A deviation that the others see only after $L$ periods earns $\Pi$ for $L + 1$ periods, so
    \omterm{def:b2:repeated-games:collusion}{collusion} requires $\beta^{L + 1} \ge 1 - 1/n$ (\cref{exo:b2:repeated-games:12}). Markets in which prices are public and
    frequent, such as posted prices, make \omterm{def:b2:repeated-games:collusion}{collusion} easier than markets with secret, negotiated contracts.
  \item \emph{Demand growth.} If demand grows by a factor $\gamma$ each period, future collusive profits are worth more
    relative to today's gain, and the condition becomes $\beta\gamma \ge 1 - 1/n$ (\cref{exo:b2:repeated-games:7}).
  \item \emph{Quantity competition.} With Cournot competition the punishment, reversion to the Cournot
    equilibrium, is milder than price competition at cost, and so is the deviator's gain. For two firms with
    linear demand the critical value is $9/17 = 0.529$ (\cref{exo:b2:repeated-games:11}).
\end{itemize}

\begin{omfigure}
\begin{tikzpicture}[econ]
\begin{axis}[omrecord, ybar, bar width=7pt, xmin=-0.6, xmax=7.6, ymin=0, ymax=1.15, width=0.7\linewidth,
    height=0.4\linewidth, xlabel={Period}, ylabel={Profit (share of $\Pi$)}, xtick={0,...,7},
    ytick={0,0.5,1}, ymajorgrids=false, legend style={at={(0.5,-0.3)}, anchor=north, legend columns=2,
    /tikz/every even column/.append style={column sep=10pt}}]
  \addplot[fill=omDef!55, draw=omDef] coordinates {(0,0.5) (1,0.5) (2,0.5) (3,1) (4,0) (5,0) (6,0) (7,0)};
  \addplot[fill=omThm!45, draw=omThm] coordinates {(0,0.5) (1,0.5) (2,0.5) (3,1) (4,1) (5,0) (6,0) (7,0)};
  \draw[guide] (axis cs:-0.6,0.5) -- (axis cs:7.6,0.5);
  \node[anchor=south west, font=\footnotesize] at (axis cs:4.5,0.52) {collusive path};
  \legend{Detected at once, Detected after one period}
\end{axis}
\end{tikzpicture}

{\small Profit of a firm that deviates in period 3 from \omterm{def:b2:repeated-games:collusion}{collusion} in a duopoly competing on price. It earns half
the monopoly profit while colluding, the whole of it in the period it undercuts, and nothing once the rival
reverts to pricing at cost. A detection lag of one period adds a second period of full profit.}
\end{omfigure}

\begin{model}{Collusion in a repeated game}
\assumes Firms that discount future profits, observe one another's prices or outputs (perhaps with a lag),
and punish a deviation by reverting to competition.
\predicts \omterm{def:b2:repeated-games:collusion}{Collusion} is easier with fewer firms, more patient firms, faster detection, growing demand and more
transparent prices.
\tested Records of discovered cartels: durations of a few years on average, with a wide spread (Levenstein
and Suslow 2006).
\breaks Multiplicity: competition is also an equilibrium whatever the \omterm{def:b2:intertemporal-choice:discount-factor}{discount factor}, so the model says when
\omterm{def:b2:repeated-games:collusion}{collusion} \emph{can} be sustained, not whether it is. Cartels break up less often from cheating than from entry
and from failing to adjust the agreement when conditions change (Levenstein and Suslow 2006).
\end{model}

\begin{record}[How long cartels last]
% ledger: b2:ls2006-t1, b2:ls2006-text
Levenstein and Suslow (2006) review studies of cartel duration. Across them the average is about five years. Many
cartels break up within a year, many others last five to ten, and some last decades. In their own sample of 42
international cartels prosecuted by the United States or the European Commission in the 1990s, the average was
5 years with a range of 0 to 20. A price-fixing agreement in aluminium phosphide lasted from January to November
1990 and collapsed when a major producer refused to take part; cartels in cartonboard (1986--91), graphite
electrodes (1992--97) and sorbates (1979--97) lasted years. The authors find that cheating and poor monitoring
end some cartels, but that entry and the difficulty of adjusting an agreement to changed conditions are the
larger threats.
\end{record}

\section{Reputation}\label{sec:b2:repeated-games:reputation}

\begin{definition}[Reputation]\label{def:b2:repeated-games:reputation}
A player's \emph{reputation}\index{reputation} is what others believe about her future behaviour, or about
her type, on the basis of her past actions.
\end{definition}

Repetition matters also when the same seller meets a new buyer each time. A buyer cannot see the quality of a
good before buying, and a seller saves a cost by delivering low quality. Played once, the seller cheats and
buyers, foreseeing it, do not buy. Suppose instead that every buyer can see the seller's record and buys only from
sellers who have never cheated. A seller who earns a price of 100 dalens and spends 30 on honest delivery earns 70 a
period by staying honest; cheating once earns 100 and then nothing. Honesty pays when $70 \ge (1 - \beta) \times 100$, that is, $\beta
\ge 0.3$. The buyers punish, though each buys only once, because the record makes the seller's past visible to all
of them.

When buyers are also unsure what kind of seller they face --- careful or careless, honest by nature or by
calculation --- the record is evidence, and buyers update their beliefs on it. A first complaint against a seller
with a clean record is then news, while a fifth is not.

\begin{record}[Ratings on online markets]
% ledger: b2:ch2006
Cabral and Horta\c{c}su (2006) followed the feedback records of sellers on eBay, an online auction site, from
October 2002 to March 2003, for four goods: notebook computers, two kinds of collectors' coins and a soft toy.
Negative ratings were rare: the median seller had one, and on average they were 0.9 per cent of comments. When a
seller received a first negative rating, the weekly growth of his sales fell from about $+7$ per cent to about $-7$
per cent; later negatives arrived 25 per cent faster than the first and had much less effect, and sellers with
worse records were more likely to leave the market.
\end{record}

The pattern fits a market in which the record carries information: the first negative tells buyers most, and a
seller whose \omterm{def:b2:repeated-games:reputation}{reputation} has fallen has less to lose by letting quality slip. A rating system, run by the market's
owner, is a private answer to the problem of hidden quality, without a regulator.

\section{Exercises}\label{sec:b2:repeated-games:exercises}

\begin{exercise}[$\star$]\label{exo:b2:repeated-games:1}
In a prisoner's dilemma with $T = 6$, $R = 4$, $P = 1$, $S = 0$, find the \omterm{def:b2:repeated-games:critical-delta}{critical discount factor} for grim trigger.
\end{exercise}

\begin{exercise}[$\star$]\label{exo:b2:repeated-games:2}
Read the opening chart. Which sample has the shortest and which the longest average duration? What is the longest
duration of a single cartel in the samples that report a range?
\end{exercise}

\begin{exercise}[$\star$]\label{exo:b2:repeated-games:3}
Predict, then compute: the \omterm{def:b2:repeated-games:critical-delta}{critical discount factor} for \omterm{def:b2:repeated-games:collusion}{collusion} among five firms competing on price; and for
two firms when a deviation is detected only after three periods.
\end{exercise}

\begin{exercise}[$\star$]\label{exo:b2:repeated-games:4}
The model predicts that \omterm{def:b2:repeated-games:collusion}{collusion} is easier among fewer firms. What evidence on discovered cartels would count
against it, and why might the evidence be hard to read?
\end{exercise}

\begin{exercise}[$\star\star$]\label{exo:b2:repeated-games:5}
Show by \omterm{def:b2:dynamic-games:backward-induction}{backward induction} that in the prisoner's dilemma repeated three times, both players defect in every period
of every \omterm{def:b2:dynamic-games:spe}{subgame-perfect equilibrium}. Which step fails if neither player knows when the game ends?
\end{exercise}

\begin{exercise}[$\star\star$]\label{exo:b2:repeated-games:6}
Both players use \omterm{def:b2:repeated-games:tit-for-tat}{tit for tat}, with $T = 5$, $R = 3$, $P = 1$, $S = 0$ and $\beta = 0.6$. A player defects in period 0 and then
plays \omterm{def:b2:repeated-games:tit-for-tat}{tit for tat}. Find the sequence of her stage payoffs and its \omterm{def:b2:repeated-games:average-payoff}{discounted average payoff}. Does the deviation pay?
For which $\beta$ does it not? Compare with grim trigger at $\beta = 0.6$.
\end{exercise}

\begin{exercise}[$\star\star$]\label{exo:b2:repeated-games:7}
Two firms compete on price, and demand grows by a factor $\gamma = 1.2$ a period. Derive the condition for \omterm{def:b2:repeated-games:collusion}{collusion},
and say whether firms with $\beta = 0.45$ can collude with this growth and without it.
\end{exercise}

\begin{exercise}[$\star\star$]\label{exo:b2:repeated-games:8}
In the eBay record, the notebook sellers' average weekly growth of sales was $7.12$ per cent in the four weeks before the
first negative rating and $-6.76$ per cent in the four weeks after, both as log differences. By what factor did sales
change over each period of four weeks? What does a model of \omterm{def:b2:repeated-games:reputation}{reputation} with uncertain types predict about the effect
of later negatives?
\end{exercise}

\begin{exercise}[$\star\star$]\label{exo:b2:repeated-games:9}
A cartel survives from one year to the next with probability 0.9 for reasons outside the firms' control, and the interest
rate is 5 per cent. Find the \omterm{def:b2:intertemporal-choice:discount-factor}{discount factor} relevant for \omterm{def:b2:repeated-games:collusion}{collusion}, and the largest number of firms competing on price
that can collude.
\end{exercise}

\begin{exercise}[$\star\star\star$]\label{exo:b2:repeated-games:10}
Show that in every \omterm{def:b2:static-games:nash-equilibrium}{Nash equilibrium} of the repeated prisoner's dilemma, finite or infinite, each player's \omterm{def:b2:repeated-games:average-payoff}{discounted
average payoff} is at least $P$.
\end{exercise}

\begin{exercise}[$\star\star\star$]\label{exo:b2:repeated-games:11}
Two Cournot firms face demand $p = a - Q$ and marginal cost $c$. They collude by each producing half the monopoly output,
and revert to the Cournot equilibrium for ever after a deviation. Show that the \omterm{def:b2:repeated-games:critical-delta}{critical discount factor} is $9/17$,
whatever $a$ and $c$.
\end{exercise}

\begin{exercise}[$\star\star\star$]\label{exo:b2:repeated-games:12}
Derive the condition $\beta^{L + 1} \ge 1 - 1/n$ for \omterm{def:b2:repeated-games:collusion}{collusion} among $n$ firms competing on price when a deviation is detected after
$L$ periods. What is the longest lag at which two firms with $\beta = 0.95$ per month can collude?
\end{exercise}

\section{Problem: two cement makers of Brenmoor}\label{sec:b2:repeated-games:problem}

\begin{problem}[{Weekend problem --- two cement makers supply Brenmoor, the capital of Brenland, an invented economy:
can they hold output down without an agreement, how does a slow flow of information change that, and what happens
when a third firm arrives}]
\label{pb:b2:repeated-games:1}
Demand is $p = 200 - Q$, with $p$ in dalens per tonne and $Q$ in thousands of tonnes a year; each firm's unit cost is 40 dalens.
Firms choose outputs once a year. The \omterm{def:b2:intertemporal-choice:discount-factor}{discount factor} is 0.8 a year, covering both interest and the chance that the
market changes. Profits are in thousands of dalens a year.

\medskip
\noindent\textbf{Part I --- The stage game.}
\begin{enumerate}
  \item Find the Cournot equilibrium: each firm's output, the price and each firm's profit.
  \item Find the outputs that maximise joint profit if split equally, the price and each firm's profit.
  \item If one firm produces its collusive output, what is the other's best output, the price and its profit?
  \item What does the firm that kept to the collusive output earn in that case?
  \item Is the collusive outcome a \omterm{def:b2:static-games:nash-equilibrium}{Nash equilibrium} of the \omterm{def:b2:repeated-games:stage-game}{stage game}?
  \item Explain why the \omterm{def:b2:repeated-games:stage-game}{stage game} is a prisoner's dilemma.
\end{enumerate}

\medskip
\noindent\textbf{Part II --- Trigger strategies.} Each firm produces its collusive output as long as both have done so,
and its Cournot output for ever after any deviation.
\begin{enumerate}[resume]
  \item Write the condition for \omterm{def:b2:repeated-games:collusion}{collusion} to be an equilibrium.
  \item Find the \omterm{def:b2:repeated-games:critical-delta}{critical discount factor}.
  \item Can the two firms collude with a \omterm{def:b2:intertemporal-choice:discount-factor}{discount factor} of 0.8?
  \item Compare the \omterm{def:b2:repeated-games:average-payoff}{discounted average payoff} of colluding with that of the best deviation.
  \item Why might the \omterm{def:b2:intertemporal-choice:discount-factor}{discount factor} be lower than interest rates alone suggest?
  \item Is the threat to revert to Cournot output credible?
\end{enumerate}

\medskip
\noindent\textbf{Part III --- A detection lag.} Outputs are not observed directly; a deviation is noticed only after $L$
years, from falling prices and lost customers.
\begin{enumerate}[resume]
  \item Write the condition for \omterm{def:b2:repeated-games:collusion}{collusion} with a lag of $L$ years.
  \item Can the firms collude at the monopoly price with $L = 1$?
  \item With $L = 2$?
  \item With $L = 2$, find the highest price the firms can sustain with equal outputs.
  \item What is the longest lag that sustains the monopoly price?
  \item Why does a public price list help the firms?
\end{enumerate}

\medskip
\noindent\textbf{Part IV --- A third firm.} A third cement maker with the same cost enters; outputs are again observed at
once.
\begin{enumerate}[resume]
  \item Find the Cournot equilibrium with three firms.
  \item Find each firm's output and profit under \omterm{def:b2:repeated-games:collusion}{collusion}.
  \item Find the best deviation and its profit.
  \item Find the \omterm{def:b2:repeated-games:critical-delta}{critical discount factor}.
  \item Can the three firms collude at the monopoly price with $\beta = 0.8$ and no lag? With a lag of one year? Of two?
  \item State the result: the \omterm{def:b2:repeated-games:critical-delta}{critical discount factor} with a lag of $L$ years, and the highest price sustainable with
    two firms and a lag of two years.
\end{enumerate}
\end{problem}
